% Options for packages loaded elsewhere
\PassOptionsToPackage{unicode}{hyperref}
\PassOptionsToPackage{hyphens}{url}
\PassOptionsToPackage{dvipsnames,svgnames,x11names}{xcolor}
\documentclass[
  spanish,
  10pt,
  a4paper,
]{article}
\usepackage{xcolor}
\usepackage[top=22mm,bottom=22mm,left=22mm,right=22mm]{geometry}
\usepackage{amsmath,amssymb}
\setcounter{secnumdepth}{-\maxdimen} % remove section numbering
\usepackage{iftex}
\ifPDFTeX
  \usepackage[T1]{fontenc}
  \usepackage[utf8]{inputenc}
  \usepackage{textcomp} % provide euro and other symbols
\else % if luatex or xetex
  \usepackage{unicode-math} % this also loads fontspec
  \defaultfontfeatures{Scale=MatchLowercase}
  \defaultfontfeatures[\rmfamily]{Ligatures=TeX,Scale=1}
\fi
\usepackage{lmodern}
\ifPDFTeX\else
  % xetex/luatex font selection
\fi
% Use upquote if available, for straight quotes in verbatim environments
\IfFileExists{upquote.sty}{\usepackage{upquote}}{}
\IfFileExists{microtype.sty}{% use microtype if available
  \usepackage[]{microtype}
  \UseMicrotypeSet[protrusion]{basicmath} % disable protrusion for tt fonts
}{}
\usepackage{setspace}
\makeatletter
\@ifundefined{KOMAClassName}{% if non-KOMA class
  \IfFileExists{parskip.sty}{%
    \usepackage{parskip}
  }{% else
    \setlength{\parindent}{0pt}
    \setlength{\parskip}{6pt plus 2pt minus 1pt}}
}{% if KOMA class
  \KOMAoptions{parskip=half}}
\makeatother
\usepackage{longtable,booktabs,array}
\newcounter{none} % for unnumbered tables
\usepackage{calc} % for calculating minipage widths
% Correct order of tables after \paragraph or \subparagraph
\usepackage{etoolbox}
\makeatletter
\patchcmd\longtable{\par}{\if@noskipsec\mbox{}\fi\par}{}{}
\makeatother
% Allow footnotes in longtable head/foot
\IfFileExists{footnotehyper.sty}{\usepackage{footnotehyper}}{\usepackage{footnote}}
\makesavenoteenv{longtable}
\ifLuaTeX
\usepackage[bidi=basic,shorthands=off]{babel}
\else
\usepackage[bidi=default,shorthands=off]{babel}
\fi
\ifLuaTeX
  \usepackage{selnolig} % disable illegal ligatures
\fi
\setlength{\emergencystretch}{3em} % prevent overfull lines
\providecommand{\tightlist}{%
  \setlength{\itemsep}{0pt}\setlength{\parskip}{0pt}}

\usepackage{amsmath,amssymb,mathtools}
\usepackage{microtype}
\usepackage{fancyhdr}
\usepackage{needspace}
\usepackage{etoolbox}
\usepackage{float}
\usepackage{caption}
\captionsetup{font=small,labelformat=empty,skip=7pt}
\floatplacement{figure}{H}
\definecolor{DocBlue}{HTML}{173F57}
\definecolor{DocGray}{HTML}{52616C}
\pagestyle{fancy}
\fancyhf{}
\fancyhead[L]{\small\sffamily\color{DocGray} pySNSPD | Desarrollo del modelo}
\fancyhead[R]{\small\sffamily\color{DocGray} Revisión 0.3}
\fancyfoot[L]{\small\sffamily\color{DocGray} 8 de septiembre de 2026}
\fancyfoot[R]{\small\sffamily\thepage}
\setlength{\headheight}{15pt}
\setlength{\emergencystretch}{2em}
\setlength{\parskip}{5pt plus 1pt}
\setlength{\parindent}{0pt}
\widowpenalty=10000
\clubpenalty=10000
\displaywidowpenalty=10000
\predisplaypenalty=10000
\renewcommand{\arraystretch}{1.16}
\AtBeginEnvironment{longtable}{\small}
\pretocmd{\section}{\Needspace{6\baselineskip}}{}{}
\pretocmd{\subsection}{\Needspace{5\baselineskip}}{}{}
\AtBeginDocument{\hypersetup{colorlinks=true,linkcolor=DocBlue,urlcolor=DocBlue}}
\allowdisplaybreaks[1]
\usepackage{bookmark}
\IfFileExists{xurl.sty}{\usepackage{xurl}}{} % add URL line breaks if available
\urlstyle{same}
\hypersetup{
  pdftitle={D. Síntesis{,} plan y verificaciones},
  pdflang={es},
  colorlinks=true,
  linkcolor={Maroon},
  filecolor={Maroon},
  citecolor={Blue},
  urlcolor={Blue},
  pdfcreator={LaTeX via pandoc}}

\title{D. Síntesis, plan y verificaciones}
\usepackage{etoolbox}
\makeatletter
\providecommand{\subtitle}[1]{% add subtitle to \maketitle
  \apptocmd{\@title}{\par {\large #1 \par}}{}{}
}
\makeatother
\subtitle{Desarrollo del modelo pySNSPD · Documento 4 de 5}
\author{}
\date{8 de septiembre de 2026 · Revisión 0.3}

\begin{document}
\maketitle

{
\setcounter{tocdepth}{2}
\tableofcontents
}
\setstretch{1.07}
\newpage

\section{D.0. Qué cambia y para
qué}\label{d.0.-quuxe9-cambia-y-para-quuxe9}

La revisión 0.3 organiza la propuesta en tres tareas físicas: \textbf{A
determina el espectro y la energía estática; B describe las ocupaciones
y los intercambios de energía; C acopla el condensado, la corriente y el
circuito.} El objetivo sigue siendo mejorar el modelamiento de SNSPD. La
construcción microscópica es una herramienta para evaluar y ampliar los
cierres del detector; su utilidad final debe medirse en los observables.

Esta iteración añade un quinto documento, E, que permite trabajar los
conceptos necesarios mediante figuras, ejemplos resueltos y ejercicios
con respuesta editable. Sus cinco temas comienzan abiertos. El avance
del cuaderno y la validación física del modelo son registros distintos:
resolver un ejercicio no valida un transitorio, y obtener un residuo
pequeño no demuestra que un concepto se haya comprendido.

{\def\LTcaptype{none} % do not increment counter
\begin{longtable}[]{@{}
  >{\raggedright\arraybackslash}p{(\linewidth - 4\tabcolsep) * \real{0.3333}}
  >{\raggedright\arraybackslash}p{(\linewidth - 4\tabcolsep) * \real{0.3333}}
  >{\raggedright\arraybackslash}p{(\linewidth - 4\tabcolsep) * \real{0.3333}}@{}}
\toprule\noalign{}
\begin{minipage}[b]{\linewidth}\raggedright
Documento
\end{minipage} & \begin{minipage}[b]{\linewidth}\raggedright
Cambio principal
\end{minipage} & \begin{minipage}[b]{\linewidth}\raggedright
Pregunta de lectura
\end{minipage} \\
\midrule\noalign{}
\endhead
\bottomrule\noalign{}
\endlastfoot
A & Derivación de la referencia normal; alcance estático; significado de
\(\lambda\) y del ajuste de Simon & ¿Qué energía se varía y qué
información microscópica contiene? \\
B & Notación FD; burbuja, retención y Fano; traslado de la energía
adiabática & ¿Qué se distribuye, qué se conserva y qué puede escapar? \\
C & Ecuación gTDGL de partida y forma estable; conexión explícita con la
fuerza propuesta & ¿Qué parte se conserva y qué cierre se propone
ampliar? \\
D & Mapa de cambios, evidencia y próximos pasos & ¿Qué afirmaciones
están verificadas y con qué alcance? \\
E & Cinco unidades pedagógicas con evaluación pendiente & ¿Se puede
explicar y usar cada concepto sin confundir sus objetos? \\
\end{longtable}
}

\section{D.1. Acuerdos físicos de esta
versión}\label{d.1.-acuerdos-fuxedsicos-de-esta-versiuxf3n}

\subsection{D.1.1. Una corriente justificada y una ampliación por
comprobar}\label{d.1.1.-una-corriente-justificada-y-una-ampliaciuxf3n-por-comprobar}

La corriente Usadel incorporada al cierre modificado tiene una
justificación de continuidad. Vodolazov desarrolla el término correctivo
en la página 11 de {[}V{]}; la memoria explica la cancelación
correspondiente en su anexo C {[}M{]}. La revisión 0.3 conserva esa
justificación. Comparar corrientes críticas de aproximaciones distintas
no basta para declarar inválida una de ellas.

El funcional común añade una pregunta: si se calculan la fuerza de
amplitud y la corriente desde una misma energía, ¿mejora la consistencia
del almacenamiento y su acoplamiento con las ocupaciones? A construye la
parte estática a partir de {[}F{]}. C plantea cómo usarla sin confundir
esa construcción con una derivación microscópica completa de la
movilidad temporal.

En una rama uniforme, suave y térmica, las identidades que se quieren
respetar son

\[
\begin{gathered}
X_{|\Delta|}^{\rm FD}
=\left.\frac{\partial f_e^{\rm FD}}{\partial|\Delta|}\right|_{T,\mathbf q},
\qquad
\left.\frac{\partial f_e^{\rm FD}}{\partial\mathbf q}\right|_{T,|\Delta|}
=\frac{\hbar}{2e}\mathbf j_s,\\
\frac{\partial X_{|\Delta|}^{\rm FD}}{\partial q}
=\frac{\hbar}{2e}\frac{\partial j_s}{\partial|\Delta|}.
\end{gathered}\tag{D.1}
\]

La segunda línea es igualdad de derivadas mixtas, a dirección de
corriente fija. No sustituye la ecuación espacial de conservación de
carga. A su vez, ni esta identidad ni la continuidad fijan por sí solas
la disipación del condensado.

\subsection{D.1.2. Tres objetos que la notación debe
separar}\label{d.1.2.-tres-objetos-que-la-notaciuxf3n-debe-separar}

{\def\LTcaptype{none} % do not increment counter
\begin{longtable}[]{@{}
  >{\raggedright\arraybackslash}p{(\linewidth - 4\tabcolsep) * \real{0.3333}}
  >{\raggedright\arraybackslash}p{(\linewidth - 4\tabcolsep) * \real{0.3333}}
  >{\raggedright\arraybackslash}p{(\linewidth - 4\tabcolsep) * \real{0.3333}}@{}}
\toprule\noalign{}
\begin{minipage}[b]{\linewidth}\raggedright
Símbolo
\end{minipage} & \begin{minipage}[b]{\linewidth}\raggedright
Objeto
\end{minipage} & \begin{minipage}[b]{\linewidth}\raggedright
Qué permite calcular
\end{minipage} \\
\midrule\noalign{}
\endhead
\bottomrule\noalign{}
\endlastfoot
\(f_{\rm FD}(E,T)\) & Probabilidad de ocupación, adimensional &
Poblaciones térmicas y factores de bloqueo \\
\(f_e^{\rm FD}\) & Densidad de energía libre electrónica & Fuerzas
termodinámicas y entropía \\
\(u_e\) & Densidad de energía interna electrónica & Almacenamiento y
balance de energía \\
\end{longtable}
}

A deriva la contribución normal antes de añadir la diferencia
superconductora. B incorpora la extensión adiabática fuera de
equilibrio:

\[
\begin{gathered}
u_e=U_{\rm vac}(|\Delta|,q)+u_{\rm qp},\\
u_{\rm qp}=4N_0\int_0^\infty E\rho(E;|\Delta|,q)f(E)\,dE.
\end{gathered}\tag{D.2}
\]

El factor cuatro corresponde a la DOS normal por espín y al conteo de
excitaciones positivas de A--B. El término de vacío y el de excitaciones
son partes de una misma contabilidad electrónica. Los fonones y el
circuito exterior tienen sus propios almacenamientos. La dependencia de
\(U_{\rm vac}\) con el superflujo no debe volver a añadirse como otra
energía local idéntica.

Una temperatura equivalente es una etiqueta de energía a espectro
fijado. Si \(u_{\rm qp}=\mathcal U(T_E;|\Delta|,q)\) y
\(\partial_T\mathcal U>0\), se puede invertir ese mapa. \textbf{La
inversión no recupera una distribución no térmica completa.} B.5 muestra
estados de igual energía con fuerzas QP distintas; C.4 explica cómo
avanzar variables energéticas y consultar una temperatura cuando hace
falta un coeficiente térmico.

\subsection{D.1.3. Material, cascada y condensado: dónde entra cada
dato}\label{d.1.3.-material-cascada-y-condensado-duxf3nde-entra-cada-dato}

La DOS fonónica cuenta modos. El espectro \(\alpha^2F\) incorpora además
el peso de interacción electrón--fonón; su integral ponderada define el
acoplamiento adimensional \(\lambda\). A.5 y E.5 explican esta
diferencia. El apéndice de Simon consultado en A reescala resultados de
acoplamiento débil por \(2.1/1.76\) para ajustar el gap; no resuelve el
sistema completo de Eliashberg con funciones dependientes de frecuencia
{[}S{]}. Este ajuste resulta útil como referencia material, pero no
permite atribuirle por sí solo todas las propiedades de un funcional de
acoplamiento fuerte.

La condición inicial de burbuja resume una cascada previa. B.3 y B.8
separan absorción, reparto entre sectores y escape al sustrato, usando
los alcances concretos de Vodolazov y Allmaras {[}V, AT, AP{]}. En un
sistema electrónico y fonónico aislado puede fluctuar el número de QP
mientras la energía total retenida permanece fija. Por ello una varianza
de conteo Fano no se transforma automáticamente en una varianza de la
energía total inyectada.

Un depósito compacto tampoco demuestra que cada evento tenga un único
máximo espacial. La figura B.5 compara perfiles construidos de igual
integral para enseñar esa diferencia. No calcula una tasa de eventos con
dos lóbulos ni la distancia entre ellos en NbN. Una predicción de ese
tipo requiere un modelo de cascada y transporte con geometría e
interacciones especificadas.

\section{D.2. Mapa de lectura y traslado de
ecuaciones}\label{d.2.-mapa-de-lectura-y-traslado-de-ecuaciones}

La reorganización evita que A deba enseñar a la vez espectro, cinética y
dinámica. Se conservan las etiquetas de A que no se trasladan; los
saltos de numeración son deliberados para facilitar comparar revisiones.

{\def\LTcaptype{none} % do not increment counter
\begin{longtable}[]{@{}
  >{\raggedright\arraybackslash}p{(\linewidth - 4\tabcolsep) * \real{0.3333}}
  >{\raggedright\arraybackslash}p{(\linewidth - 4\tabcolsep) * \real{0.3333}}
  >{\raggedright\arraybackslash}p{(\linewidth - 4\tabcolsep) * \real{0.3333}}@{}}
\toprule\noalign{}
\begin{minipage}[b]{\linewidth}\raggedright
Ubicación en 0.2
\end{minipage} & \begin{minipage}[b]{\linewidth}\raggedright
Ubicación en 0.3
\end{minipage} & \begin{minipage}[b]{\linewidth}\raggedright
Contenido
\end{minipage} \\
\midrule\noalign{}
\endhead
\bottomrule\noalign{}
\endlastfoot
A.1--A.14 & A.1--A.14 & Funcional estático y derivadas; se añaden
A.7a--A.7c para la referencia normal \\
Sección A.4; A.15--A.24 & Sección B.11; B.43--B.52 & Energía adiabática,
coordenadas espectrales y fuerzas a ocupaciones fijadas \\
A.16a & B.44a & Identidad auxiliar del traslado anterior \\
Sección A.5; A.25--A.31 & Sección C.12; C.35--C.41 & Corriente,
continuación uniforme y comparación de cierres \\
A.25a--A.25c & C.35a--C.35c & Desarrollo de fase y continuidad \\
Sección A.6; A.32--A.34 & Sección A.5; A.32--A.34 & Material,
\(\alpha^2F\), \(\lambda\) y acoplamiento fuerte \\
Figura A.2 & Figura C.4 & Comparación uniforme de corriente \\
Apertura de C & C.0a--C.0f & gTDGL de partida, forma estable y puente
hacia C.9 \\
\end{longtable}
}

La ecuación final consultada en la memoria está en el \textbf{anexo C,
ecuación (C.3), página impresa 121, página 152 del PDF}. Los
coeficientes usados en C proceden de sus ecuaciones (2.29)--(2.31),
página impresa 31. Se indica la localización efectiva de la copia
consultada para evitar ambigüedad al buscar anexos. C coteja también la
ecuación (3.24) de Allmaras, página impresa 94 {[}AT{]}.

La forma estable de C conserva el término temporal
\(\partial_t|\Delta|^2\,\Delta\) y la multiplicación del lado derecho
por su factor de relajación. La propuesta se conecta después con la
fuerza de amplitud. La equivalencia algebraica de esa escritura y la
continuidad no certifican cualquier discretización cerca de
\(|\Delta|=0\); ese tratamiento sigue siendo un punto específico de
implementación.

Para estudiar el conjunto: E.2--E.3 prepara A; E.4--E.5 prepara B;
después C permite seguir el acoplamiento al detector. E.1 aporta el
contexto de campos, partículas y excitaciones. También puede leerse A--C
de forma continua y consultar E sólo cuando aparezca una dificultad
concreta.

\section{D.3. Qué se ejecutó en esta
revisión}\label{d.3.-quuxe9-se-ejecutuxf3-en-esta-revisiuxf3n}

Se reejecutaron en Geminga los tres cálculos ligeros heredados de 0.2 y
se añadieron las figuras y comprobaciones del cuaderno. Se utilizó el
entorno científico existente, con Python 3.10.20, NumPy 2.2.6, SciPy
1.15.3 y Matplotlib 3.10.8, limitando los hilos de cálculo a uno. La
carpeta de trabajo es \nolinkurl{sandbox/model_v0_3}, dentro de
\nolinkurl{/home/jdiaz/pysnspd}.

{\def\LTcaptype{none} % do not increment counter
\begin{longtable}[]{@{}
  >{\raggedright\arraybackslash}p{(\linewidth - 4\tabcolsep) * \real{0.3333}}
  >{\raggedright\arraybackslash}p{(\linewidth - 4\tabcolsep) * \real{0.3333}}
  >{\raggedright\arraybackslash}p{(\linewidth - 4\tabcolsep) * \real{0.3333}}@{}}
\toprule\noalign{}
\begin{minipage}[b]{\linewidth}\raggedright
Comprobación
\end{minipage} & \begin{minipage}[b]{\linewidth}\raggedright
Resultado registrado en 0.3
\end{minipage} & \begin{minipage}[b]{\linewidth}\raggedright
Lo que verifica
\end{minipage} \\
\midrule\noalign{}
\endhead
\bottomrule\noalign{}
\endlastfoot
A: derivadas de la energía uniforme & Error relativo máximo
\(1.34\times10^{-10}\) en tres puntos & Fuerza y corriente frente a
derivadas numéricas, en esos puntos \\
A/C: corriente máxima Usadel & \(38.8503\,\mu\)A; estable entre 400 y
3200 términos Matsubara & Convergencia del problema uniforme de
acoplamiento débil con los parámetros declarados \\
B: 3000 eventos de cada reacción & Residuo relativo de intercambio
\(3.96\times10^{-20}\) & Signos y conservación en eventos sintéticos \\
B: invertir energía térmica & Error absoluto máximo
\(2.22\times10^{-16}\) en unidades de la prueba & Inversión a espectro
fijo y capacidad positiva en el intervalo ensayado \\
B: estados de igual energía & Razón de fuerzas QP \(25.1873\) entre dos
paquetes elegidos & Energía sola no fija todos los momentos de la
distribución \\
C: celda aislada ilustrativa & Deriva de energía \(6.75\times10^{-13}\);
diferencia entre coordenadas \(1.14\times10^{-12}\) o menor, en unidades
adimensionales & Balance y equivalencia numérica del modelo de celda
declarado \\
C: circuito con resistencia prescrita & Residuo integrado de balance
\(1.67\times10^{-5}\) o menor, en unidades adimensionales & Contabilidad
de potencia para esa entrada temporal y cuadratura \\
E: ejemplos resueltos & BdG: \(E=\pm5\), pesos \(0.8/0.2\); curvatura
reducida \(1\); autovalores de rigidez adimensionales \(1,3\) & Álgebra
de los ejemplos, sin resolver las actividades evaluadas \\
\end{longtable}
}

Los números coincidentes con 0.2 son \textbf{reproducción de pruebas
anteriores}, no nuevas predicciones materiales. En la prueba instantánea
de B, el residuo de la regla de la cadena es \(6.94\times10^{-18}\); su
pequeñez se interpreta a precisión numérica, no como cero matemático
inferido de una ejecución.

Los parámetros de la comparación uniforme están en
\nolinkurl{verificaciones/A_resumen.json}; B y C declaran sus unidades y
modelos en \nolinkurl{B_verificaciones.json} y \nolinkurl{C_checks.json}.
Las tablas CSV permiten rehacer los gráficos. La celda polinómica, el
circuito prescrito y los paquetes sintéticos se identifican como tales
en los pies de figura. No se ejecutó un transitorio espacial del
detector, una cascada estocástica ni un cálculo DFT/DFPT nuevo.

\subsection{D.3.1. Figuras y controles de
entrega}\label{d.3.1.-figuras-y-controles-de-entrega}

El conjunto incorpora 17 figuras distintas: dos en A, cinco en B, cuatro
en C y seis en E. Siete son recursos nuevos para B/E; la antigua
comparación A.2 acompaña ahora C.12. Cada figura se conserva en PNG y
PDF vectorial. Los cinco documentos tienen fuentes Markdown editables y
LaTeX compilable.

La comprobación de entrega verifica que todas las etiquetas de ecuación
aparezcan en los PDF, que las figuras estén incrustadas y que no haya
avisos de símbolos ausentes o desbordes. La revisión visual y el conteo
final de páginas se registran en
\nolinkurl{verificaciones/QA_estructura.json} y
\nolinkurl{verificaciones/QA_visual.json}. El manifiesto de la entrega
identifica los archivos y sus hashes; el paquete conserva la estructura
del repositorio para que las rutas relativas de las imágenes sigan
funcionando.

\section{D.4. Próximos pasos del
modelo}\label{d.4.-pruxf3ximos-pasos-del-modelo}

\begin{enumerate}
\def\labelenumi{\arabic{enumi}.}
\tightlist
\item
  \textbf{Cerrar las convenciones estáticas.} Contrastar la derivación
  de A con la elección material: normalización de DOS, gap, rama
  espectral y régimen de corriente. Decidir después si el primer
  catálogo utiliza acoplamiento débil, un ajuste fenomenológico
  identificado o una construcción de acoplamiento fuerte más completa.
\item
  \textbf{Elegir el estado retenido de B.} Comparar una distribución
  energética resuelta con cierres por momentos bajo las mismas entradas
  materiales. Medir errores en potencia electrón--fonón y fuerza de
  amplitud, además del error en energía. Una temperatura equivalente
  puede seguir usándose para lectura o evaluación de coeficientes.
\item
  \textbf{Probar el acoplamiento local de C.} Sustituir la celda
  pedagógica por un cierre microscópico explícito, verificar balance
  aislado y disipación, y precisar la movilidad. Mantener una prueba de
  equivalencia de la escritura temporal antes de introducir cambios
  espaciales.
\item
  \textbf{Integrar espacio y circuito.} Revisar continuidad, fronteras,
  paso temporal, extensión cerca de ceros del condensado y contabilidad
  de energía con el circuito. Sólo entonces comparar transitorios,
  latencias y señales con la referencia anterior y con datos.
\item
  \textbf{Evaluar la condición inicial cuando haya un modelo de
  cascada.} Separar retención media, fluctuaciones de reparto, escape y
  forma espacial. La hipótesis de varios lóbulos debe contrastarse como
  distribución de eventos; no basta cambiar el ancho de una única
  burbuja.
\end{enumerate}

Estos pasos son propuestas de desarrollo, no cambios realizados en el
solver. La revisión local de código utilizada como contexto es
\nolinkurl{391bb301a6c36f4438754a51a10a7fe2aade0555}. Los módulos de
producción no forman parte de esta actualización documental.

\section{D.5. Seguimiento del
aprendizaje}\label{d.5.-seguimiento-del-aprendizaje}

E dispone de cinco temas con identificadores estables E01 a E05. Cada
uno incluye propósito, figura pedagógica, desarrollo, ejemplo resuelto,
tres actividades y criterio de comprensión. Las respuestas se escriben
en su Markdown, debajo de cada enunciado.

{\def\LTcaptype{none} % do not increment counter
\begin{longtable}[]{@{}
  >{\raggedright\arraybackslash}p{(\linewidth - 4\tabcolsep) * \real{0.3333}}
  >{\raggedright\arraybackslash}p{(\linewidth - 4\tabcolsep) * \real{0.3333}}
  >{\raggedright\arraybackslash}p{(\linewidth - 4\tabcolsep) * \real{0.3333}}@{}}
\toprule\noalign{}
\begin{minipage}[b]{\linewidth}\raggedright
Tema
\end{minipage} & \begin{minipage}[b]{\linewidth}\raggedright
Estado de 0.3
\end{minipage} & \begin{minipage}[b]{\linewidth}\raggedright
Evidencia pendiente
\end{minipage} \\
\midrule\noalign{}
\endhead
\bottomrule\noalign{}
\endlastfoot
E01. Modelo Estándar & ABIERTO & Respuestas y explicación de campos
frente a ocupaciones \\
E02. Nambu y espín & ABIERTO & Cálculos y distinción entre índices y
estados físicos \\
E03. Cálculo funcional & ABIERTO & Variación, bordes y uso correcto de
la envolvente \\
E04. Electrón y hueco & ABIERTO & Contabilidad de carga y energía;
vínculo con superconductividad \\
E05. DFPT y DOS & ABIERTO & Modos, normalización espectral y pesos de
interacción \\
\end{longtable}
}

El estado cambia a \textbf{CERRADO / APROBADO} con al menos 8 de 10
puntos y el criterio esencial del tema satisfecho. En la siguiente
revisión se añadirá la pauta correspondiente a las respuestas recibidas,
con corrección razonada y nueva oportunidad de resolver lo pendiente. Un
tema sin respuestas permanece abierto y sin nota. Si se detecta después
una confusión esencial, el registro permite reabrirlo indicando el
motivo; el historial no se borra.

\section{D.6. Fuentes y trazabilidad}\label{d.6.-fuentes-y-trazabilidad}

{[}M{]} \emph{Multiscale Modeling of the Transient Response of
Superconducting Nanowire Single-Photon Detectors}, copia
\nolinkurl{memoria_02.pdf} de \nolinkurl{/home/jdiaz/memoria/main}, 173
páginas. Se consultaron las ecuaciones (2.29)--(2.31), (C.3) y el
desarrollo de continuidad del anexo C. Las páginas impresas y las
páginas del PDF se distinguen en C.

{[}V{]} D. Yu. Vodolazov, \emph{Single-Photon Detection by a Dirty
Current-Carrying Superconducting Strip Based on the Kinetic-Equation
Approach}, Phys. Rev.~Applied \textbf{7}, 034014 (2017).
\href{https://arxiv.org/pdf/1611.06060}{Preprint consultado}: ecuaciones
(16)--(22) para condiciones iniciales y página 11, ecuación (36), para
el término correctivo.

{[}F{]} P. Virtanen, A. Vargunin y M. Silaev, \emph{Quasiclassical free
energy of superconductors: Disorder-driven first-order phase transition
in superconductor/ferromagnetic-insulator bilayers}, Phys. Rev.~B
\textbf{101}, 094507 (2020), ecuación (20).
\href{https://doi.org/10.1103/PhysRevB.101.094507}{Artículo} y
\href{https://arxiv.org/pdf/1909.00992v1}{preprint v1}, donde el título
es \emph{Quasiclassical expressions for the free energy of
superconducting systems}. A explicita la reducción utilizada.

{[}AT{]} J. P. Allmaras, \emph{Modeling and Development of
Superconducting Nanowire Single-Photon Detectors}, tesis doctoral,
Caltech (2020).
\href{https://thesis.caltech.edu/13748/08/Allmaras_Thesis_Final.pdf}{PDF
oficial}. B localiza las secciones de cascada, escape y condición
inicial; C utiliza la ecuación (3.24) y su discusión numérica.

{[}AP{]} J. P. Allmaras et al., \emph{Intrinsic Timing Jitter and
Latency in Superconducting Nanowire Single-Photon Detectors}, Phys.
Rev.~Applied \textbf{11}, 034062 (2019).
\href{https://doi.org/10.1103/PhysRevApplied.11.034062}{Artículo} y
\href{https://arxiv.org/pdf/1805.00130v2}{preprint v2}. Las referencias
específicas de fluctuaciones Fano y escape están desarrolladas en B.

{[}S{]} Simon et al., \href{https://arxiv.org/pdf/2501.13791v3}{preprint
arXiv:2501.13791v3}, sección VII.3, ecuaciones (11)--(14). La numeración
citada corresponde a esta versión del preprint. A distingue el ajuste
del gap de una solución completa de Eliashberg.

Los documentos 0.1 y 0.2, las observaciones recibidas y los PDF de las
fuentes se conservan como procedencia. El paquete 0.3 distribuye los
documentos y figuras originales de esta revisión; no redistribuye los
PDF de terceros. Todos los archivos de la entrega y las comprobaciones
se copian a Geminga bajo las mismas carpetas de versión.

\end{document}
